%! AP Statistics — Unit 1, Lesson 1.2 — Student Packet Cover Sheet
\documentclass[10pt]{article}
\usepackage{apstats-article}
\usepackage{apstats-boxes}

%=============================================================================
\begin{document}

% ── Full-bleed navy banner ────────────────────────────────────────────────────
\begin{tikzpicture}[remember picture, overlay]
  \fill[navy] (current page.north west)
    rectangle ([yshift=-0.9in]current page.north east);
\end{tikzpicture}
\vspace{-0.8in}

\begin{center}
  {\color{white}\textbf{\LARGE AP Statistics}}\\[4pt]
  {\color{skymid}\large\textbf{Unit 1: Exploring One-Variable Data}}\\[6pt]
  {\color{white}\textbf{\large Lesson 1.2 \quad The Language of Variation --- Variables}}
\end{center}

\vspace{0.22in}

\namedateperiod

\vspace{0.18in}

% ── LEARNING TARGETS ─────────────────────────────────────────────────────────
\begin{learningtargetbox}
\small
\begin{enumerate}[leftmargin=1.5em, itemsep=5pt]
  \item \ldots{} classify a variable as \textit{categorical} or \textit{quantitative}
        and explain my reasoning using the ``average test.''
  \item \ldots{} distinguish between \textit{discrete} and \textit{continuous}
        quantitative variables and give an example of each.
  \item \ldots{} explain why variable type must be identified before choosing a
        statistical display or summary.
\end{enumerate}
\end{learningtargetbox}

\vspace{0.14in}

% ── TABLE OF CONTENTS ────────────────────────────────────────────────────────
\begin{tocbox}
\small
\renewcommand{\arraystretch}{1.5}
\begin{tabularx}{\linewidth}{@{} c l X r @{}}
\rowcolor{sky}
\textbf{\#} & \textbf{Component} & \textbf{Description} & \textbf{Score} \\
\midrule
1 & Warm-Up        & Spiral review + prior knowledge connection    & \blank{1.2cm} \\
2 & Guided Notes   & Direct instruction: variable types \& examples & \blank{1.2cm} \\
3 & Group Activity & Card sort + mini-whiteboard round              & \blank{1.2cm} \\
4 & Exit Ticket    & Individual classification with justification   & \blank{1.2cm} \\
5 & Homework       & NBA dataset analysis + AP-style practice       & \blank{1.2cm} \\
\midrule
  & \multicolumn{2}{r}{\textbf{Total}} & \blank{1.2cm} \\
\end{tabularx}
\end{tocbox}

\vspace{0.14in}

% ── KEEP IN MIND ─────────────────────────────────────────────────────────────
\begin{remindbox}
\small
Variable type is the \emph{first} question in data analysis. Get it wrong and every
display, summary, and conclusion that follows will be wrong too.

\vspace{6pt}
\begin{tabularx}{\linewidth}{@{} >{\centering\arraybackslash\columncolor{goldbg}}p{0.06\linewidth}
                                    X @{}}
\textbf{\large!} &
\textbf{Numbers are not always quantitative.}\enspace
ZIP codes, jersey numbers, and student IDs are numbers that label --- averaging
them produces a meaningless result. \\[6pt]
\textbf{\large!} &
\textbf{Use the average test.}\enspace
Ask: ``Does the average of these values tell me something useful?''
If yes, the variable is quantitative. If no, it is categorical. \\[6pt]
\textbf{\large!} &
\textbf{Discrete vs.\ continuous matters later.}\enspace
In Unit 4 (Probability), discrete variables get probability histograms;
continuous variables get density curves. Build the habit now. \\[6pt]
\textbf{\large!} &
\textbf{Variable type determines your tools.}\enspace
Categorical $\to$ bar chart, two-way table, proportion.\enspace
Quantitative $\to$ histogram, boxplot, mean, standard deviation. \\
\end{tabularx}
\end{remindbox}

\end{document}
